\scanpage{110}
\uline{Corollaire 1}\quad Supposons que $X$, \inserted{$Y$} \deleted{\unsure{soient des}} \inserted{soient} des espaces $T_0$, \deleted{\unsure{élémentaires qui}} admettant des bases d'ouverts formées d'ens \deleted{quasi-cpts} rétrocompacts et q. cpts pour la top. constructible [donc compacts pour la dite], \deleted{\unsure{lesquels}} \inserted{Conditions} équivalentes:

a) $f:X\longrightarrow Y$ \deleted{soit} \deleted{\unsure{une telle que}} \inserted{est} quasi-compacte

\deleted{[\unsure{si} cpt $X$}\quad \deleted{Alors} b) $f^{\mathrm{cons}}$ est une \deleted{application} application propre d'espaces loc. cpts \deleted{c) \unsure{les} $f(\unsure{X_i})$ \unsure{parties} $X\subset f$}.

\deleted{[\unsure{les} $f$ transforme parties \unsure{antifiques|antifinies} en parties \unsure{antifiques|antifinies}]}

\margin{c) $f$ transforme parties \unsure{antifiques|antifinies} en parties \unsure{antifiques|antifinies}, et les fibres de $f$ sont quasi-compactes [\unsure{après}: Bourbaki, Top \deleted{\unsure{II}} I, \unsure{10,2,th.1}]}

[NB. $f$ en particulier transforme les parties constructibles en parties constructibles, \deleted{et \unsure{bien les cas fermés} \inserted{\unsure{plus loc}} \unsure{est fermé}, \unsure{cela est transforme les}}]

\uline{Corollaire 2}\quad Sous les conditions précédentes \inserted{\unsure{les vérifier}}, supposons $Y$ et $X$ quasi cpts, donc $X^{\mathrm{cons}}$, $Y^{\mathrm{cons}}$ \uline{compacts}. Conditions équivalentes

a) $f$ transforme parties constructibles en parties constructibles [i.e. transforme ``ouverts et fermés'' en ``ouverts et fermés'']

b) $f$ transforme parties constructibles en parties réunions de constructibles [i.e. ``ouverts et fermés'' en ``ouverts'']

(c) $f^{\mathrm{cons}}$ est une application \uline{ouverte}. \deleted{\unsure{(qui est propre)}}

Ces conditions sont vérifiées si $f$ est \uline{bijective}.

\uline{Cor. 3}\quad Si les conditions du lemme 1 sont remplies, alors (i) $f(X)$ est \unsure{antifique|antifinie}; (ii) pour \unsure{$E\subset Y$}, une partie \inserted{$E$} de $f(X)$ est \unsure{antifique|antifinie} ssi $f^{-1}(E)$ l'est (iii) Si $f$ est surjective, une partie de $Y$ est constructible ssi $f^{-1}(\unsure{Y|E})$ l'est.

\margin{\figure{\begin{tikzpicture}[x=1cm,y=1cm]
\draw(0,0)--(.12,-5.0);
\draw(0,-6.0)..controls(.1,-7.8)and(-.3,-11.0)..(.4,-11.3);
\end{tikzpicture}}}
\margin{109}
